% chapters/requirements/rf_permissions.tex
% RF-5: Permisos y autorización
% ============================================================================

\item \textbf{Permisos y autorización.} \label{rf:perms}
El sistema de autorización opera en tres capas contextuales, eliminando el concepto de
<<rol global>> rígido. Los permisos se derivan de: (1)~el \textit{flag}
\texttt{is\_staff}, que otorga permisos de gestor a nivel de organización;
(2)~la \texttt{SubjectMembership}, que define el rol (\texttt{COORDINATOR},
\texttt{TEACHER}, \texttt{STUDENT}) y los permisos contextuales dentro de una asignatura
concreta mediante un campo \acs{json} \texttt{permission\_overrides}; y (3)~las reglas de
asignación de corrección (\texttt{AssignmentRule}), que restringen el acceso a instancias
o problemas concretos. Cada petición se evalúa contra el contexto actual (organización,
curso, asignatura) para determinar si el usuario tiene autorización. La
\cref{tab:perms_defaults} recoge los permisos contextuales y sus valores por defecto según
el rol.

\begin{table}{tab:perms_defaults}{Permisos contextuales y valores por defecto según rol}
    \begin{tabular}{lccc}
        \hline
        \textbf{Permiso} & \textbf{COORD} & \textbf{TEACHER} & \textbf{STUDENT} \\
        \hline
        \texttt{can\_create\_exams}       & Sí & No & No \\
        \texttt{can\_manage\_models}      & Sí & No & No \\
        \texttt{can\_manage\_call\_list}  & Sí & No & No \\
        \texttt{can\_assign\_correctors}  & Sí & No & No \\
        \texttt{can\_grade}               & Sí & Sí & No \\
        \texttt{can\_annotate}            & Sí & Sí & No \\
        \texttt{can\_resolve\_issues}     & Sí & No & No \\
        \texttt{can\_publish\_grades}     & Sí & No & No \\
        \texttt{can\_manage\_reviews}     & Sí & No & No \\
        \texttt{can\_view\_all\_instances} & Sí & No & No \\
        \texttt{can\_export\_grades}      & Sí & No & No \\
        \texttt{can\_manage\_members}     & Sí & No & No \\
        \texttt{can\_delegate\_perms}     & Sí & No & No \\
        \hline
    \end{tabular}
\end{table}

\begin{functional}

    % RF-5.1 ----------------------------------------------------------------
    \item \textbf{Permisos de gestor (\texttt{is\_staff}).} \label{rf:perms:staff}
    El sistema debe otorgar automáticamente permisos de administración a los usuarios con
    \texttt{is\_staff=True}.

    \textbf{Entrada:} evento interno disparado al establecer \texttt{is\_staff=True} en un
    usuario.

    \textbf{Proceso:} un usuario con \texttt{is\_staff=True} obtiene permisos para
    gestionar usuarios, cursos, configuración y auditoría dentro de su organización. Tiene
    acceso de lectura y escritura a todas las asignaturas de su organización sin necesidad
    de una \texttt{SubjectMembership} explícita.

    \textbf{Salida:} el usuario puede realizar operaciones de gestor.

    % RF-5.2 ----------------------------------------------------------------
    \item \textbf{Asignación de permisos contextuales por defecto.}
    \label{rf:perms:defaults}
    El sistema debe asignar automáticamente un conjunto de permisos por defecto al crear una
    \texttt{SubjectMembership}, basado en el rol.

    \textbf{Entrada:} creación de una \texttt{SubjectMembership} con un rol
    (\texttt{COORDINATOR}, \texttt{TEACHER} o \texttt{STUDENT}).

    \textbf{Proceso:} el sistema inicializa el campo \texttt{permission\_overrides} como un
    diccionario vacío. Internamente, una función
    \texttt{get\_effective\_permissions(membership)} combina los permisos base del rol
    (véase \cref{tab:perms_defaults}) con las anulaciones en
    \texttt{permission\_overrides}. Los permisos base por defecto de cada rol se obtienen de
    la configuración de la organización (véase \cref{rf:config:default_perms}), lo que
    permite a cada organización adaptar los valores iniciales. Los permisos del rol
    \texttt{STUDENT} son fijos y no admiten \texttt{permission\_overrides} individuales;
    las restricciones de acceso del estudiante se configuran a nivel de examen (véase
    \cref{rf:exams:student_perms}).

    \textbf{Salida:} la \texttt{SubjectMembership} queda creada con los permisos efectivos
    correspondientes.

    % RF-5.3 ----------------------------------------------------------------
    \item \textbf{Modificación de permisos contextuales.}
    \label{rf:perms:override}
    El sistema debe permitir al coordinador de una asignatura modificar los permisos de un
    profesor en esa asignatura, añadiendo o eliminando permisos específicos.

    \textbf{Entrada:} petición PATCH al recurso de la \texttt{SubjectMembership} con un
    objeto \acs{json} de permisos a modificar (por ejemplo,
    \texttt{\{can\_create\_exams: true\}}).

    \textbf{Proceso:} el sistema verifica que el usuario que modifica es el coordinador de
    la asignatura. Fusiona el \acs{json} recibido con el
    \texttt{permission\_overrides} existente. Solo se pueden otorgar permisos que el
    otorgante posee (véase \cref{rf:perms:delegation}). Se registra el cambio en el
    \textit{log} de auditoría.

    \textbf{Salida:} respuesta con la lista actualizada de permisos efectivos del profesor
    en esa asignatura.

    % RF-5.4 ----------------------------------------------------------------
    \item \textbf{Revocación de permisos contextuales.} \label{rf:perms:revoke}
    El sistema debe permitir al coordinador revocar permisos contextuales previamente
    otorgados a un profesor, restaurando los valores por defecto.

    \textbf{Entrada:} petición PATCH al recurso de la \texttt{SubjectMembership} con los
    permisos a revocar (por ejemplo, \texttt{\{can\_create\_exams: null\}}). Un valor
    \texttt{null} elimina la anulación y restaura el permiso por defecto del rol.

    \textbf{Proceso:} el sistema elimina las claves indicadas del campo
    \texttt{permission\_overrides}.

    \textbf{Salida:} respuesta con la lista actualizada de permisos efectivos.

    % RF-5.5 ----------------------------------------------------------------
    \item \textbf{Consulta de permisos efectivos.} \label{rf:perms:query}
    El sistema debe permitir consultar los permisos efectivos de un usuario en una
    asignatura concreta.

    \textbf{Entrada:} petición GET al recurso de permisos de la \texttt{SubjectMembership}.

    \textbf{Proceso:} el sistema calcula los permisos efectivos combinando: el \textit{flag}
    \texttt{is\_staff}, los permisos base del rol en \texttt{SubjectMembership} y las
    anulaciones en \texttt{permission\_overrides}.

    \textbf{Salida:} respuesta con la lista de permisos efectivos del usuario, indicando la
    fuente de cada uno (base o anulación).

    % RF-5.6 ----------------------------------------------------------------
    \item \textbf{Restricción de delegación.} \label{rf:perms:delegation}
    El sistema debe impedir que un usuario otorgue permisos que no posee.

    \textbf{Entrada:} cualquier petición de modificación de \texttt{permission\_overrides}.

    \textbf{Proceso:} antes de aplicar la asignación, el sistema compara la lista de
    permisos solicitados con los permisos efectivos del otorgante. Si algún permiso
    solicitado no está en los permisos del otorgante, la operación se rechaza íntegramente.

    \textbf{Salida:} en caso de intento de exceso, respuesta de error con código
    \acs{http}~403.

    % RF-5.7 ----------------------------------------------------------------
    \item \textbf{Permisos de corrección (capa \texttt{AssignmentRule}).}
    \label{rf:perms:correction}
    El sistema debe restringir el acceso a las instancias de examen según las reglas de
    asignación de corrección.

    \textbf{Entrada:} cualquier petición de acceso a una instancia de examen (consulta de
    \acs{pdf}, anotaciones, calificaciones).

    \textbf{Proceso:} el sistema verifica que el usuario que accede es el coordinador, un
    gestor o un corrector asignado a esa instancia o problema (mediante
    \texttt{AssignmentRule}). Los estudiantes solo acceden a su propia instancia y solo
    durante los períodos permitidos (publicación, revisión), limitados a los permisos
    configurados en el examen. Un profesor que no es corrector de una instancia puede
    consultar sus atributos pero no el \acs{pdf}.

    \textbf{Salida:} acceso concedido si se cumplen las condiciones; en caso contrario,
    respuesta de error con código \acs{http}~403.

\end{functional}
